% !PRINTED 132
\weberpage{132}
If, as before, $h$ is the degree of $v$, and if we set $\lambda = h$ when $h$ is divisible by $p$, and $\lambda = hp$ when $h$ is not divisible by $p$, then it follows from (7) that $(va)^\lambda = 1$, i.e. the degree of $va$ is a divisor of $h$ or of $hp$, hence in any case a divisor of $p\nu$; i.e. $va$ \gesperrt{is contained in} $V$.

If we now replace $a$ in (7) by $a^x$ and let $x$ run through the sequence of numbers $0, 1, \ldots p-1$, then $a^x$ runs through the group $Q$ and $va^x$ through the system $vQ$, and for every $\lambda$ we have
\weq{8}{(va^x)^\lambda = v^\lambda a^{x\lambda}.}
Now if the degree $h$ of $v$ is not divisible by $p$, then
\[
(va^x)^h = a^{xh},
\]
and hence $= 1$ only when $x = 0$; on the other hand
\[
(va^x)^{ph} = 1,
\]
i.e. the degree of $va^x$, when $x$ is not $= 0$, is a divisor of $ph$, but not of $h$, hence divisible by $p$; on the other hand it is $= h$ when $x = 0$.

But if the degree $h$ of $v$ is divisible by $p$, then it is of the form $pg$, and $g$ is not divisible by $p$, because $p$ divides $n$, and hence also $h$, only once. Thus $v^g$ is an element of $R$ of degree $p$, and must therefore, by \S.\,30, be contained in $Q$. If accordingly we set
\[
v^g = a^y,
\]
then
\[
(va^x)^g = a^{y+gx},
\]
and this is $= 1$ if and only if $y + gx \equiv 0 \ (\mathrm{mod}\ p)$, which occurs for only \gesperrt{one} value of $x$. For this value of $x$ the degree of $va^x$ is a divisor of $g$; for the other values of $x$ it is a divisor of $pg$, but not of $g$, hence divisible by $p$. With this the first two assertions of theorem $\beta$) are proved.

In order to see the last part as well, one may observe that if $v_1, v_2$ denote two elements of $V$, the systems $v_1 Q, v_2 Q$ are either wholly identical, or have not a single element in common. From this it follows that $V$ can be represented in the following manner:
\[
V = v_1 Q + v_2 Q + \cdots + v_{r_2} Q
\]
(although $V$ is in general not a group), and in each of these $r_2$ partial systems $vQ$ there are $p-1$ elements whose degree
%CONT>>
